\subsection{Simplicial four-polytopes: an explicit contractible family}
\label{sec:gy1985-stacking}

This development belongs to item~1985-658.10 and problem~GP-1985-42.
The notebook on P001230 says only that even simplicial polytopes in
\(\mathbb R^4\) can have \ru{более или менее любой} topology of their
fixed-type realization spaces. The qualification means ``more or less any'';
it does not specify a homeomorphism, homotopy equivalence, stable equivalence,
or the normalization of realizations. An editorial homotopy formulation is:
does every nonempty finite simplicial-complex homotopy type occur as a
normalized realization space of a simplicial four-polytope? The stronger
precise conjecture of \citet[Conjecture~4]{adiprasitopadrol2017universality}
asks for stable equivalence to every nonempty open primary basic
semialgebraic set defined over \(\mathbb Z\), even requiring the four-polytope
to be neighborly. This latter formulation is a literature conjecture, not
an exact quotation from the seminar. Adiprasito--Padrol use a homogenized
realization model: vector configurations realizing the cone over the
polytope, modulo \(\mathrm{GL}_{d+1}(\mathbb R)\). The fixed affine-frame
and prescribed-history model proved contractible below is specified
separately; no identification of these normalization conventions is
assumed.

\paragraph{The established results and their limits.}
Mn\"ev's 1985 announcement already gives configuration universality and a
simplicial-polytope statement with dimension unrestricted
\citep[Theorems~3--4]{mnev1985varieties}. Adiprasito--Padrol explicitly regard
the earlier arguments for the simplicial part as incomplete and supply
the missing proof; the 1985 announcement must therefore not be treated
as equivalent to an independently checked complete proof of that part. The proof of general four-polytope
universality reported in \citet[Main Theorem]{richtergebertziegler1995realization}
concerns work of 1994 and appeared in 1995; it does not impose simpliciality.
\citet[Theorem~2]{adiprasitopadrol2017universality} proves the neighborly
simplicial result in dimension \(2n-4\) with \(2n\) vertices. That paper first
circulated in February 2014, appeared online in May 2016, and was printed
in 2017. Its Conjecture~4 preserves the difficult fixed dimension.
Their discussion after Conjecture~4 also cites established disconnected
realization spaces of simplicial four-polytopes, so nontrivial topology
in dimension four is already known; universality is the remaining
quantifier. The same paper records trivial realization spaces for stacked polytopes
in the discussion after Conjecture~4; its Theorem~3 proves a ball result
for specified sewing families. The proof below supplies an explicit
continuous insertion and contraction for a narrowly specified stacking
space. It is not a new universality theorem or a new discovery of
contractibility for stacked polytopes.

As of November 2026, no resolution of fixed-dimensional simplicial
universality was found in the bounded search for this chapter. The latest
opened source explicitly retaining the constant-dimension simplicial
question is the 2024 compendium, entries CG68--CG69
\citep{schaefercardinalmiltzow2024existential}. This is dated evidence, not a
claim to have excluded every subsequent publication. A later 2026 search
lead concerning noncontractible realization spaces is recorded as
unverified in the research log and is not used in the argument.

\paragraph{The space to which the proof applies.}
Fix \(d\geq2\), and fix the labelled simplex
\(P_0=\operatorname{conv}(0,e_1,\ldots,e_d)\subset\mathbb R^d\).
A \emph{prescribed stacking history} \(\mathcal H\) of length \(r\) specifies,
at step \(k\), a labelled simplex facet \(F_k\) of the combinatorial polytope
produced by the previous steps. The next labelled point must lie strictly
beyond \(F_k\) and strictly beneath every other current facet.
Let \(\mathcal R_0\) be the one-point space consisting of the fixed simplex.
Recursively, \(\mathcal R_k\) consists of the coordinates of the first \(k\)
inserted points satisfying every one of these conditions at its own step.
These are \emph{history-compatible} realizations. Coordinates of the initial
simplex remain fixed, so no affine group factor is present. We do not
identify \(\mathcal R_r\) with the full realization space of the final
combinatorial type without a separate deletion argument.

For a configuration \(q\in\mathcal R_{k-1}\), write \(P(q)\) for its convex
hull. For each facet \(G\), choose an inward affine supporting form
\[
  \ell_{G,q}(x)=a_G(q)\cdot x+b_G(q),\qquad
  \|a_G(q)\|=1,
\]
which vanishes on \(G\) and is positive on the interior of \(P(q)\).
These choices are continuous in \(q\). Indeed, the labelled vertices of a
simplex facet give a nonzero affine cofactor vector. Divide its spatial
part by its norm and choose its sign to be positive at the average of all
vertices of \(P(q)\). That average is in the interior; the sign test never
vanishes, so it is continuous, including if the parameter space were
initially disconnected. This construction avoids an unproved global
choice of a facet normal.

\begin{lemma}[The stacking insertion region \ours]
For a full-dimensional simplicial polytope \(P\) and a facet \(F\), set
\[
 U_F(P)=\{x\in\mathbb R^d:\ell_F(x)<0,
                  \ \ell_G(x)>0\text{ for every facet }G\ne F\}.
\]
This is a nonempty open convex set. For every \(x\in U_F(P)\),
\(\operatorname{conv}(P\cup\{x\})\) has the stacking combinatorial type:
it retains all facets except \(F\), and replaces \(F\) by the simplices
\(\operatorname{conv}(R\cup\{x\})\), one for each ridge \(R\) of \(F\).
\end{lemma}
\begin{proof}
The inequalities define an intersection of finitely many open half-spaces,
hence an open convex set. Let \(b\) be the barycenter of the vertices of
\(F\), and let \(n_F=-a_F\) be the outward unit normal. Since \(b\) lies in
the relative interior of \(F\), \(\ell_G(b)>0\) for every other facet.
For sufficiently small \(\varepsilon>0\), the point
\(b+\varepsilon n_F\) therefore satisfies all the required inequalities;
nonemptiness is established quantitatively in the next proposition.

For the combinatorial assertion, a segment from \(x\) to any point of \(P\)
meets the plane \(\ell_F=0\) at a point of \(F\), unless it is wholly on the
\(P\) side. At the crossing, all other facet inequalities hold because they
hold at both endpoints. Consequently
\[
 \operatorname{conv}(P\cup\{x\})
       =P\cup\operatorname{conv}(F\cup\{x\}).
\]
Every old facet \(G\ne F\) remains a supporting facet, since
\(\ell_G(x)>0\). If a ridge \(R\subset F\) is shared by \(F\) and \(G\),
then
\[
 h_R(y)=\ell_G(x)\ell_F(y)-\ell_F(x)\ell_G(y)
\]
is nonnegative on \(P\), vanishes precisely on \(R\) there, and also
vanishes at \(x\). It exposes exactly
\(\operatorname{conv}(R\cup\{x\})\), which is a simplex facet. The boundary
of the union displayed above consists of the old boundary without the
relative interior of \(F\), and these lateral faces of the added pyramid.
Thus the listed facets exhaust the new boundary and give the stated type.
\end{proof}

\begin{proposition}[A continuous section and a fiberwise contraction \ours]
The map
\[
 \pi_k:\mathcal R_k\longrightarrow\mathcal R_{k-1}
\]
forgetting the last vertex is surjective and has a continuous section.
Moreover, \(\mathcal R_k\) strongly deformation retracts onto the graph of
that section while every previous vertex is held fixed.
\end{proposition}
\begin{proof}
Let \(F=F_k\). Its facet barycenter \(b(q)\) and outward unit normal
\(n_F(q)=-a_F(q)\) are continuous. Define
\[
 \varepsilon(q)=\frac12\min\left(
       1,\ \min_{G\ne F}
       \frac{\ell_{G,q}(b(q))}{1+|a_G(q)\cdot n_F(q)|}
                         \right)>0,
 \qquad s(q)=b(q)+\varepsilon(q)n_F(q).
\]
The minimum is over the fixed finite labelled facet set at this step, so
\(\varepsilon\) and \(s\) are continuous. For the selected facet,
\(\ell_{F,q}(s(q))=-\varepsilon(q)<0\). For every other facet, write
\(u=|a_G(q)\cdot n_F(q)|\). Then
\[
 \ell_{G,q}(s(q))
 \geq \ell_{G,q}(b(q))-\varepsilon(q)u
 \geq \ell_{G,q}(b(q))
             \left(1-\frac{u}{2(1+u)}\right)>0.
\]
Thus \(s(q)\in U_F(P(q))\), proving both surjectivity and the section
\(q\mapsto(q,s(q))\).

The fiber over \(q\) is exactly \(U_F(P(q))\). Its convexity proves that
\[
 H_k((q,x),t)=\bigl(q,(1-t)x+t\,s(q)\bigr),\qquad 0\leq t\leq1,
\]
remains in \(\mathcal R_k\). This map is jointly continuous, equals the
identity at \(t=0\), has image in the section at \(t=1\), and fixes every
point of the section for every \(t\). These are precisely the strong
deformation-retraction conditions. No claim that contractible fibers
alone imply a homotopy equivalence is needed.
\end{proof}

\begin{theorem}[Contractibility for a prescribed stacking history \ours]
For every finite prescribed facet-stacking history on the fixed simplex,
the normalized history-compatible realization space \(\mathcal R_r\) is
nonempty and contractible. In particular this holds in dimension four.
\end{theorem}
\begin{proof}
The space \(\mathcal R_0\) is a point. At each step the section proves
nonemptiness, and the preceding proposition retracts \(\mathcal R_k\)
onto a copy of \(\mathcal R_{k-1}\). If a contraction of
\(\mathcal R_{k-1}\) is already given, first use \(H_k\) to move to the
section and then apply the contraction there through the continuous
section map. Concatenating these two homotopies contracts
\(\mathcal R_k\). Induction gives a finite, continuous contraction of
\(\mathcal R_r\). The conclusion concerns exactly the recursively defined
space: no assertion that every realization of the final type admits
this history has entered the proof.
\end{proof}

\paragraph{A second check of the section estimate.}\ours\
The argument can be rederived without unit normals, which is useful for
exact arithmetic. Take inward forms with nonzero gradients \(a_G\), put
\[
 \eta=\min\left(1,\ \min_{G\ne F}
          \frac{\ell_G(b)}{2(1+|a_G\cdot a_F|)}\right),
 \qquad x=b-\eta a_F.
\]
Then \(\ell_F(x)=-\eta\|a_F\|^2<0\), while
\(\ell_G(x)\geq\ell_G(b)-\eta|a_G\cdot a_F|>0\).
This independently checks the sign and the strict-slack bound. If the
vertices and inward forms are rational, \(b,\eta,x\) are rational; no
square root is required. These identities check the proof for every
admissible step, independently of the finite examples computed elsewhere
in this chapter. Numerical or exact examples do not establish equality
with the full realization space.

\begin{proposition}[Rational points in simplicial realization spaces \ours]
Every nonempty realization space of a labelled simplicial polytope
contains a rational realization. The same holds after normalizing a
specified affinely independent vertex frame to a rational simplex, on
any nonempty such chart.
\end{proposition}
\begin{proof}
Start with a labelled simplicial realization. Each facet has exactly
\(d\) affinely independent vertices, and its supporting affine form is
strictly positive on every other vertex. These are finitely many open
conditions on the vertex coordinates: affine independence and the signs
of the appropriate determinants. They continue to hold in a sufficiently
small neighbourhood. Every listed simplex facet therefore persists.

For completeness, this persistence rules out extra facets. Every ridge
of a listed simplex facet belongs to exactly two listed facets in the
original simplicial boundary complex. In the perturbed convex hull those
two facets still meet along this ridge, and a ridge of a convex polytope
belongs to exactly two facets. Hence no facet adjacent across a ridge to
a listed facet can be unlisted. The facet-adjacency graph of a convex
polytope is connected, so every facet is listed. Each old labelled vertex
remains a vertex because it is a vertex of a persisting simplex facet.
Thus the combinatorial type is constant throughout that neighbourhood.

Rational points are dense in the coordinate space, so choose all vertex
coordinates rational within this neighbourhood. For a normalized chart,
first apply the unique affine map sending the chosen vertex frame to the
specified rational frame. Hold those frame coordinates fixed and perturb
the remaining coordinates rationally within the same finite strict
inequalities. This gives the normalized rational realization as well.
\end{proof}

\begin{corollary}[The arithmetic limit of a simplicial universality claim \ours]
Stable equivalence over \(\mathbb Q\), generated by rational equivalences
and stable projections with nonempty relatively open polyhedral fibers
defined over the base, cannot identify a nonempty simplicial-polytope
realization space with every primary basic semialgebraic set over
\(\mathbb Z\) without an openness or comparable arithmetic restriction.
\end{corollary}
\begin{proof}
Existence of a rational point is preserved by rational equivalences
whose maps and inverses are defined on the spaces. It is also preserved
in both directions by a stable projection. One direction follows by
projection. For the other, over a rational base point the fiber is a
nonempty relatively open polyhedron with rational linear equations and
strict rational inequalities. Its affine hull has a rational point and
a rational direction basis; density in that affine hull yields a rational
point satisfying all strict inequalities.

The primary basic set
\(S=\{t\in\mathbb R:t^2-2=0,\ t>0\}\) is nonempty and has no rational
point. The proposition supplies a rational point for every nonempty
simplicial realization space, so such a space cannot be stably equivalent
to \(S\) in this sense. By contrast, a nonempty open subset of Euclidean
space always has rational points. This explains why the established and
conjectural simplicial universality statements restrict to open primary
sets. It imposes no corresponding obstruction on homotopy universality:
this particular \(S\) already has the homotopy type of a point.
\end{proof}

\paragraph{The remaining obstruction and a bounded next step.}
A fixed stacking history admits a global section and convex insertion
fibers, so it cannot encode arbitrary topology. The unresolved universal
claim needs mechanisms beyond this restricted construction. In the
known approaches, recording complicated order types inside auxiliary
configurations does not by itself ensure that a fixed simplicial
four-dimensional face lattice forces those order types. The discussion
after Conjecture~4 in \citet{adiprasitopadrol2017universality} identifies
precisely this distinction using stacked vertex figures and contractions.
The present proof does not address that enforcement problem.

A further bounded run could take one explicitly labelled small final
stacked four-polytope and determine, by exact inequalities, whether every
realization in the chosen normalization permits a specified last-vertex
deletion. A proof would identify its full realization space with the
history-compatible space for that case; a counterexample would delimit
the identification. This is a concrete extension of the certificate,
not a proposal to settle fixed-dimensional universality by checking a
finite list of examples.
